% Generated by `taskgen course examples` from dlp.diffie_hellman (seed 10).
% Do not edit: change examples.yaml and regenerate.
\begin{example}[протокол Диффи--Хеллмана]\label{ex:dh-toy}
Алиса и Боб согласуют ключ по протоколу Диффи--Хеллмана с открытыми параметрами $p = 227$, $g = 147$. Секретный показатель Алисы $a = 212$, Боба $b = 41$. \par (а) Бинарным алгоритмом вычислить открытые ключи $A$, $B$ и общий ключ $K$ --- отдельно на стороне Алисы и на стороне Боба. \par (б) Найти $\ord(g)$. Сколько различных значений может принимать общий ключ при этих $p$, $g$? Что знает Ева и как она может найти $K$?

\par\noindent\textit{Решение.} Степени вычисляем бинарным алгоритмом: биты показателя слева направо, $r \gets r^2$, при бите $1$ ещё $r \gets r \cdot (\text{основание})$; все вычисления по модулю $p = 227$.\[A = 147^{212}:\quad\begin{array}{c|rrrrrrrr}\text{бит} & 1 & 1 & 0 & 1 & 0 & 1 & 0 & 0 \\ \hline r^2 &  & 44 & 59 & 76 & 131 & 136 & 29 & 160 \\ r & 147 & 112 & 59 & 49 & 131 & 16 & 29 & 160\end{array}\qquad A \equiv 160.\]\[B = 147^{41}:\quad\begin{array}{c|rrrrrr}\text{бит} & 1 & 0 & 1 & 0 & 0 & 1 \\ \hline r^2 &  & 44 & 120 & 43 & 33 & 181 \\ r & 147 & 44 & 161 & 43 & 33 & 48\end{array}\qquad B \equiv 48.\]Алиса получает $B$ и вычисляет $K_A = B^a$, Боб получает $A$ и вычисляет $K_B = A^b$:\[K_A = 48^{212}:\quad\begin{array}{c|rrrrrrrr}\text{бит} & 1 & 1 & 0 & 1 & 0 & 1 & 0 & 0 \\ \hline r^2 &  & 34 & 33 & 181 & 212 & 225 & 136 & 109 \\ r & 48 & 43 & 33 & 62 & 212 & 131 & 136 & 109\end{array}\qquad K_A \equiv 109.\]\[K_B = 160^{41}:\quad\begin{array}{c|rrrrrr}\text{бит} & 1 & 0 & 1 & 0 & 0 & 1 \\ \hline r^2 &  & 176 & 104 & 221 & 36 & 161 \\ r & 160 & 176 & 69 & 221 & 36 & 109\end{array}\qquad K_B \equiv 109.\]$K_A = K_B = g^{ab} = 109$.\par (б) Порядок $g$ (последовательное снижение): $\ord(g) \mid p - 1 = 226 = 2 \cdot 113$; $147^{226/2} = 147^{113} \equiv 1$ $\Rightarrow o = 113$; $147^{113/113} = 147^{1} \equiv 147$ $\ne 1$. Итак, $n = \ord(g) = 113$. Ключ $K = g^{ab}$ лежит в $\langle g \rangle$, поэтому он принимает не более $113$ значений. $\langle g \rangle$ --- подгруппа простого порядка $q = 113$; так и выбирают параметры на практике (но $q$ должно быть огромным, $q \approx 2^{256}$), а $A, B$ проверяют условием $A^q = 1$.\par Ева видит $p, g, A = 160, B = 48$ (но не $a, b$). Ей достаточно найти $a' = \log_g A$ (это $a \bmod n = 99$): тогда $K = B^{a'}$. Перебором --- до $n - 1 = 112$ умножений, методом Шенкса --- около $2\lceil\sqrt{n}\,\rceil = 22$ операций; для настоящих $n \approx 2^{256}$ это недостижимо.\[A = 160,\quad B = 48,\quad K = 109.\]
\end{example}
